% Fragmento ES para inserción, no documento autónomo.
% Destino: XII, ampliacion_20260922.tex, después de la prueba de
% «Realización exacta del cuarto de giro» y antes de introducir S_zeta.
% Propietarios y comprobaciones: LEER_PRIMERO.md.
\subsubsection{Restitución del calendario y acarreo ternario}
\label{ins29:xii:acarreo}

La extracción del plano de cuarto de giro admite una descripción exacta
de la información temporal conservada. Recibimos el calendario de
\(108\) posiciones y la extensión dodecafásica sobre la fase nonádica
construidos anteriormente. Para esta descripción utilizamos representantes
de base cero:
\[
 \theta=r+9\beta,\qquad 0\leq r<9,\quad 0\leq\beta<12.
\]
La posición que el registro enumera de \(1\) a \(9\) es aquí \(r+1\).
La actualización enriquecida precede a este calendario finito y conserva,
además, la ruta, las incidencias y la memoria acumulada. En el plano ya
construido, \(C_{12}^{\beta}W=WJ_0^{\beta}\); por tanto, su lectura depende
de \(b=\beta\bmod4\).

\begin{proposition}[Fibra ternaria de la lectura de cuarto de giro]
\label{ins29:xii:fibra}
El par formado por la fase nonádica \(r\) y la fase planar \(b\) determina
un cociente de \(36\) posiciones del calendario. Su ley de composición es
\[
 (r,b)(r',b')=
 \left(r+r'\bmod9,\;
 b+b'+\left\lfloor\frac{r+r'}9\right\rfloor\bmod4\right).
\]
Cada lectura posee tres preimágenes en \(\mathbb Z/108\mathbb Z\).
La componente adicional \(k=\beta\bmod3\) restituye la coordenada
dodecafásica mediante
\[
 \beta=9b+4k\pmod{12}.
\]
\end{proposition}
\begin{proof}
La suma de dos representantes da
\(\theta+\theta'=r+r'+9(\beta+\beta')\). La división euclidiana de
\(r+r'\) por \(9\) produce el acarreo de la fórmula. La lectura
\((r,b)\) equivale al residuo \(\theta\bmod36\), pues
\(\theta\equiv r+9b\pmod{36}\). Su núcleo es
\(\{0,36,72\}\), de modo que cada fibra tiene tres elementos.
Finalmente, \(9b+4k\) es congruente con \(b\) módulo \(4\) y con
\(k\) módulo \(3\); estos dos residuos determinan un único elemento
módulo \(12\).
\end{proof}

\begin{proposition}[Localización de la extensión no escindida]
\label{ins29:xii:extension}
La descomposición
\[
 \theta\longmapsto(j,z)=(\theta\bmod4,\theta\bmod27),
 \qquad \theta=81j+28z\pmod{108},
\]
identifica el calendario con
\(\mathbb Z/4\mathbb Z\times\mathbb Z/27\mathbb Z\).
La proyección nonádica es \((j,z)\mapsto z\bmod9\), y la inclusión
dodecafásica \(\beta\mapsto9\beta\) toma la forma
\[
 \beta\longmapsto(\beta\bmod4,9\beta\bmod27).
\]
La obstrucción a escindir la extensión completa queda concentrada en
\[
 0\longrightarrow\mathbb Z/3\mathbb Z
 \longrightarrow\mathbb Z/27\mathbb Z
 \longrightarrow\mathbb Z/9\mathbb Z\longrightarrow0.
\]
\end{proposition}
\begin{proof}
Los coeficientes \(81\) y \(28\) tienen residuos \((1,0)\) y \((0,1)\)
módulo \((4,27)\), respectivamente. Esto verifica la inversa y las dos
aplicaciones indicadas. La sucesión ternaria no se escinde: un producto
\(\mathbb Z/3\mathbb Z\times\mathbb Z/9\mathbb Z\) tiene exponente
\(9\), mientras \(\mathbb Z/27\mathbb Z\) tiene exponente \(27\).
En cambio, la proyección del cociente de \(36\) posiciones sobre
\(\mathbb Z/9\mathbb Z\) admite la sección homomorfa
\(n\mapsto28n\pmod{36}\). Por tanto, la reducción elimina precisamente
esta obstrucción ternaria.
\end{proof}

La distinción se observa ya en \(\theta=0,36,72\): las tres posiciones
dan \(r=b=0\), mientras \(z\) vale \(0,9,18\), respectivamente.
La restitución del calendario requiere conservar la fibra que el plano
identifica. Además, el residuo cuaternario global satisface
\(j=r+b\pmod4\); la fase planar \(b\) se transporta junto con la fase
nonádica y su acarreo.

Así queda localizado el alcance del cuarto de giro: organiza una
representación orientada del registro y, acompañado del residuo nonádico,
recupera exactamente su cociente de \(36\) posiciones. La coordenada
ternaria adicional restituye las \(108\) posiciones. La memoria de la
dinámica enriquecida conserva a su vez la prolongación entre retornos
completos del calendario. La periodicidad de una representación y la
persistencia de la trayectoria registrada pertenecen, por ello, a lecturas distintas
del mismo transporte.
